\chapter{Building a silly little Doom}

Everything so far, assembled. The finished program is
\file{course3d/lesson07\_mini\_doom.c}: a corridor maze you walk through, enemies
that come after you, a gun that shoots them, and a HUD with a minimap. About 350
lines, and you have already met every idea in it.

\tryit{Run it first, then read this chapter with the source open next to
you.}{./course3d/build.sh 7}

\section{The level is a piece of text}

\begin{lstlisting}
static const char *MAP[MAP_ROWS] = {
    "###############",
    "#S....#.......#",
    "#.....#...E...#",
    "#.###.#.#####.#",
    ...
};
\end{lstlisting}

\lstinline|#| is a wall, \lstinline|.| is floor, \lstinline|S| is where the
player starts and \lstinline|E| is an enemy. Designing a level is now typing
characters into a text editor, and you can redesign the whole game in thirty
seconds. Plenty of shipped games stored their levels exactly like this.

\section{From map cell to world position}

The map is a grid of rows and columns. The world is metres along X, Y and Z. The
bridge between them is one function, and it is worth staring at until it is
obvious:

\begin{lstlisting}
static Vector3 CellToWorld(int col, int row, float y)
{
    return (Vector3){ col*CELL, y, row*CELL };
}
\end{lstlisting}

\begin{center}
\begin{tikzpicture}[scale=0.85,every node/.style={font=\sffamily\footnotesize}]
  \fill[rlight,rounded corners=3pt] (-0.7,-4.1) rectangle (12.6,1.5);

  % grid
  \foreach \c in {0,1,2,3} \foreach \r in {0,1,2} {
    \draw[codeframe,fill=white] (\c*0.9,-\r*0.9) rectangle (\c*0.9+0.9,-\r*0.9-0.9);
  }
  \fill[rgrey!45] (0,0) rectangle (0.9,-0.9);
  \fill[rgrey!45] (0.9*3,0) rectangle (0.9*4,-0.9);
  \node at (1.8,0.75) {the text map};
  \node[text=rgrey] at (1.8,-3.2) {column, row};

  \draw[->,rorange,line width=1.2pt] (4.0,-1.2) -- (5.4,-1.2);
  \node[above,text=rorange] at (4.7,-1.15) {CellToWorld};

  % world
  \begin{scope}[xshift=6.3cm,yshift=-1.2cm]
    \draw[->,rred,line width=1.1pt] (0,0) -- (2.6,0) node[right,text=rred] {$+X$ \ (column)};
    \draw[->,rblue,line width=1.1pt] (0,0) -- (0,-2.0) node[below,text=rblue] {$+Z$ \ (row)};
    \draw[->,rgreen,line width=1.1pt] (0,0) -- (-0.9,0.9) node[above left,text=rgreen] {$+Y$ \ (height)};
    \fill[rink] (0,0) circle (1.6pt);
    \node[text=rgrey,align=left] at (2.0,-2.6) {world metres};
  \end{scope}
\end{tikzpicture}
\end{center}

\keyidea{Map \textbf{column} becomes world \textbf{X}. Map \textbf{row} becomes
world \textbf{Z}. Y is not in the map at all, because a top-down plan has no
height in it. This one mapping is the whole relationship between a 2D level
editor and a 3D world, and it runs backwards too --- that is how the minimap
turns a 3D position back into a dot.}

\section{Reading the map once, at startup}

\lstinline|ResetGame()| walks the text and turns markers into game state:

\begin{lstlisting}
for (int row = 0; row < MAP_ROWS; row++)
    for (int col = 0; col < MAP_COLS; col++)
    {
        if (MAP[row][col] == 'S')
            g->camera.position = CellToWorld(col, row, EYE_HEIGHT);
        else if (MAP[row][col] == 'E' && g->enemyCount < MAX_ENEMIES)
        {
            g->enemies[g->enemyCount].position = CellToWorld(col, row, 0.9f);
            g->enemies[g->enemyCount].alive = true;
            g->enemyCount++;
        }
    }
\end{lstlisting}

Because all the state lives in one \lstinline|Game| struct, restarting the level
is a single call to \lstinline|ResetGame(&game)| bound to the R key. Worth
copying as a habit: put the game state in a struct, write a reset function, and
you get restarts for free.

\section{The order of the frame}

The loop is one long sequence, and the order matters:

\begin{enumerate}[leftmargin=*,itemsep=2pt]
  \item \textbf{Look} --- mouse delta into yaw and pitch, pitch clamped, then
        build \lstinline|forward|, \lstinline|right| and
        \lstinline|flatForward| (chapter 12).
  \item \textbf{Move} --- keys into a direction, normalized, times speed times
        \lstinline|dt|, resolved one axis at a time against the grid
        (chapter 13).
  \item \textbf{Shoot} --- on click, fire a ray and see what it finds.
  \item \textbf{Enemies} --- walk towards the player, hurt them on contact.
  \item \textbf{Draw the world} --- inside \lstinline|BeginMode3D|.
  \item \textbf{Draw the HUD} --- after \lstinline|EndMode3D|, in screen pixels.
\end{enumerate}

Notice that all six steps read the same \lstinline|forward| vector that step 1
computed. Compute it once per frame, near the top, and pass it around. Computing
it twice with different values is how you get a gun that shoots slightly to the
left of the crosshair.

\section{Walls}

One cube per \lstinline|#| cell, with a shade that varies by position so the
corridors do not look like a single flat mass:

\begin{lstlisting}
Vector3 centre = CellToWorld(col, row, WALL_HEIGHT/2);
int shade = 88 + ((col*7 + row*13) % 26);
DrawCube(centre, CELL, WALL_HEIGHT, CELL, wall);
DrawCubeWires(centre, CELL, WALL_HEIGHT, CELL, (Color){ 30, 26, 24, 255 });
\end{lstlisting}

Note \lstinline|WALL_HEIGHT/2| for the Y: the cube's centre is at half its
height, so it sits on the floor instead of sinking into it (chapter 9).

Floor and ceiling are two \lstinline|DrawPlane| calls in different colours. Two
lines, and suddenly you are indoors.

\section{Shooting}

\begin{lstlisting}
Ray shot = { game.camera.position, forward };

// how far is the wall in front of us?
float wallDistance = 100.0f;
for (float t = 0.1f; t < 60.0f; t += 0.1f)
{
    Vector3 p = Vector3Add(shot.position, Vector3Scale(shot.direction, t));
    if (IsWall(p.x, p.z)) { wallDistance = t; break; }
}

// the nearest enemy closer than that wall
int target = -1;
float best = wallDistance;

for (int i = 0; i < game.enemyCount; i++)
{
    if (!game.enemies[i].alive) continue;
    RayCollision hit = GetRayCollisionBox(shot, EnemyBox(game.enemies[i].position));
    if (hit.hit && hit.distance < best) { best = hit.distance; target = i; }
}
\end{lstlisting}

Two things are happening. The loop marching along the ray in 10 cm steps is a
crude \emph{raycast} against the grid --- it walks forward until it is inside a
wall cell, and that is your line of sight. Seeding \lstinline|best| with that
distance is what stops you shooting through walls: an enemy only counts if it is
nearer than the wall.

\gotcha{Stepping along a ray in fixed increments can step straight over a thin
wall if the step is bigger than the wall. Ours is 10 cm against 2 m walls, so it
is safe. Proper raycasters use a grid traversal algorithm (look up ``DDA'') that
cannot miss, and that is the next thing to learn if you want to go further down
this road.}

\section{Enemies}

The simplest AI that still feels like a game: if you can see me and I am not
touching you, walk at you.

\begin{lstlisting}
Vector3 toPlayer = Vector3Subtract(game.camera.position, e->position);
toPlayer.y = 0.0f;                       // ignore height: they walk, not fly
float distance = Vector3Length(toPlayer);

if (distance < 12.0f && distance > 0.9f)
{
    Vector3 step = Vector3Scale(Vector3Normalize(toPlayer), 1.6f*dt);
    if (CanStandAt(e->position.x + step.x, e->position.z)) e->position.x += step.x;
    if (CanStandAt(e->position.x, e->position.z + step.z)) e->position.z += step.z;
}
else if (distance <= 0.9f)
{
    game.health -= 25.0f*dt;             // four seconds of contact kills you
}
\end{lstlisting}

The pattern \lstinline|normalize(target - me) * speed * dt| is how absolutely
everything chases everything else: enemies, homing missiles, a pet following you.
Learn it once, use it forever.

Their movement is resolved per axis too --- same two lines as the player --- so
they slide along walls instead of getting stuck in corners.

\section{The HUD, or where the 2D course pays off}

Everything after \lstinline|EndMode3D()| is plain screen pixels:

\begin{itemize}[leftmargin=*,itemsep=3pt]
  \item \textbf{Crosshair} --- four short \lstinline|DrawLine| calls with a gap
        in the middle, centred on \lstinline|GetScreenWidth()/2|.
  \item \textbf{The gun} --- three rectangles at the bottom of the screen, offset
        by the same \lstinline|sinf(bobTimer)| that bobs the camera, so the
        weapon sways with your steps. A muzzle flash is two circles drawn for
        0.06 seconds after a shot.
  \item \textbf{Health bar} --- a background rectangle and a second one whose
        width is \lstinline|224.0f*health/100.0f|, turning red below 30.
  \item \textbf{Minimap} --- the text map drawn as little squares, with the
        player and enemies as dots. World X divided by the cell size gives the
        map column, which gives the screen X. Chapter 9's ``the floor is XZ''
        used in reverse.
\end{itemize}

\begin{lstlisting}
int px = ox + (int)(game.camera.position.x/CELL*tile);
int py = oy + (int)(game.camera.position.z/CELL*tile);
DrawCircle(px, py, 3, SKYBLUE);
DrawLine(px, py, px + (int)(forward.x*14), py + (int)(forward.z*14), SKYBLUE);
\end{lstlisting}

That last line is the little arrow showing which way you face: the 3D forward
vector, with its Y thrown away, drawn as a 2D line. Three dimensions and two
dimensions, in the same frame, sharing one number.

\section{Things to change first}

The point of having the source is breaking it. In rough order of fun per minute:

\begin{enumerate}[leftmargin=*,itemsep=2pt]
  \item Redraw the map. Add a room, a dead end, a long corridor.
  \item Change \lstinline|WALL_HEIGHT| to 8 and walk around. Then try 1.5.
  \item Set \lstinline|fovy| to 110 and then to 40. Feel how much the lens
        changes the game.
  \item Give the enemies different speeds, or make them stop when they are far
        away and only wake when they see you.
  \item Add a pickup: a \lstinline|P| in the map, drawn as a spinning cube
        (rotate it with \lstinline|GetTime()|), that refills ammo when you touch
        it. Everything you need is in chapters 9 and 13.
  \item Add jumping: a \lstinline|velocityY|, gravity subtracted every frame, and
        the Y axis resolved like X and Z.
\end{enumerate}
